\documentclass[11pt]{article}
\usepackage{amsmath,amssymb,amsthm,hyperref}
\newtheorem{proposition}{Proposition}
\newcommand{\E}{\mathbb E}
\newcommand{\Prb}{\mathbb P}
\title{Coefficient-preserving first hits and their original-law stopping cost}
\author{Campaign 3395: Engineer continuation}
\date{November 5, 2026}
\begin{document}
\maketitle

\paragraph{The live estimate is not proved.}
The question is the small-size assertion
\[
 S_3(\rho)\le\tfrac14
 \quad\text{for every BP1 circuit of factor }1+t
 \text{ and size }s<e^{aN},                         \tag{SF1}
\]
with explicit fixed rational $0<t\le1/64$, explicit
$0<a\le2^{-136}$ and admissible $N\ge N_0\ge2^{146}$.
No such triple is obtained here. A global FH16 modulus is not
required, and the small-size hypothesis is not a restriction on the
allowed circuit model.

The attempted first-hit proof below preserves coefficients and makes
its probability accounting explicit. Its remaining step is a signed
residual domination, not a missing implementation or an unavailable
tool. Neither the algebraic identities nor the bounded computations
prove SF1. FH12 and FH13 remain unchanged. No novelty or publishable
significance is claimed.

\paragraph{Reused results.}
Use the reviewed BP1--BP3, C2--C6, A4/A7, DS8--DS9,
SC10--SC11, FH1/FH6/FH9/FH15, GH4--GH14, and GW4--GW12.
Their certificates are not rerun. The finite two-block factors below
are the already reviewed AT12 factors, not a new small approximator.
All new assertions are proved directly; no external theorem is added.

\section{A first hit counts a coefficient once, even with sharing}
Expand the original finite arithmetic DAG into decorated parse
occurrences $\tau$. Addition chooses one child; multiplication chooses
an ordered pair of child occurrences. Keep distinct occurrences even
when their monomials agree. Write their coefficients and monomials as
$c_\tau>0,m_\tau$. This expansion is a mathematical indexing set,
not a claimed efficient construction.

Fix distinct original arithmetic gates $v_1,\ldots,v_r$ before the
sampler. Let $i(\tau)$ be the least index of a gate occurring anywhere
in $\tau$, or $\infty$ when none occurs. Repeated uses of the same
gate do not create further first hits.

\begin{proposition}
For the successive zero substitutions of GW7,
\[
 Z_i=Q^{[i-1]}-Q^{[i]}
       =\sum_{\tau:i(\tau)=i}c_\tau m_\tau.          \tag{SF2}
\]
The actual compulsory-cut expansion can be refined into these same
original occurrences without altering a coefficient. Consequently
the cut index and the first-hit index give a joint partition, not
two independent copies of the output coefficients.
\end{proposition}
\begin{proof}
A parse survives a zero substitution exactly when it contains none
of the substituted gates. Its first deletion is therefore $i(\tau)$,
which proves SF2 coefficientwise, including arbitrary powers and
nonnegative real constants.

For the cut at an original gate, distinguish its first occurrence
in the left-to-right traversal of a parse. At a multiplication gate,
the two terms of C5 partition parses into those whose left child
contains the cut gate and those whose left child does not but whose
right child does. These classes are disjoint. The unchanged and
zero-substituted child polynomials in C5 implement exactly these
classes. Addition preserves the partition. Earlier cuts remove
their entire classes first. This proves the claimed refinement
and explains why ordinary derivatives, which count every occurrence,
are not used.
\end{proof}

\section{The surviving density factor}
Let $\omega=(H,I_\rho)$ have exactly its original law, write
$U=U_I(H)$, and refine each actual cut pair into the occurrences
above. Every output occurrence is full-group. Call it bad if its
monomial contains distinct variables from the same source group.
For a bad occurrence belonging to actual cut $j$, let
$\chi_\tau(\omega)$ be its DS9 completion-survival indicator.
It retains both optional-source directions.
Let $\ell_j(\omega)$ be the original dependent height-three secant.
In particular, it is not recomputed after splitting a cut.

For the occurrences of cut $j$ and first hit $v$, define measures by
\[
 d\nu_{jv}(\omega,\tau)=
  \frac{c_\tau m_\tau(U)}{F(U)}\,d\Prb(\omega),\qquad
 d\mu_{jv}=F(U)\ell_j(\omega)\chi_\tau(\omega)\,d\nu_{jv}.
                                                               \tag{SF3}
\]
Zero-support terms are zero; all boxes contain the original one-hot
assignment, so their denominator is positive.

\begin{proposition}
For a list covering the output parses, these are exact identities:
\[
 \sum_{j,v}\int d\nu_{jv}
       =\E\frac{D_Q(U)}{F(U)}\le b_Q,\qquad
 S_3(\rho)=\sum_{j,v}\int
                \sigma(H)F(U)\ell_j\chi_\tau\,d\nu_{jv}.          \tag{SF4}
\]
For a noncovering list include the $\infty$ class. At each fixed box,
the bad first-hit charge is at most $Z_v(U)/F(U)$, and these
normalized deletion charges have total at most $3$ under BP1.
\end{proposition}
\begin{proof}
SF2 partitions all original coefficients. Compatibility filtering
is linear on this partition, so its bad part sums to $D_Q$, by A7.
Any completed bad occurrence has all its variables in $U$, hence
$\chi_\tau\le m_\tau(U)$. Its original one-hot value is zero.
A compatible full-group occurrence survives the completion exactly
when it survives the original assignment: its compulsory choices
already fix every optional choice. Thus the raw error consists
exactly of the completed bad occurrences. Multiplying their sum
by the common original $\ell_j$ gives the clipped error.
This proves SF3--SF4. Positivity and GW8 give the last assertion.
\end{proof}

This construction does not remove $F(U)=16^{|I|}$ from SF3.
The compulsory-anchor fraction in the active-bit average gives
exactly the existing $4^{|I|}$ unsigned bound, hence FH6, not a
constant coefficient. Nor can the sign cancel inside one fixed
occurrence: SC10--SC11 fixes its compulsory anchor and sign.
Cancellation must still relate different occurrences.

In particular retain the exact signed quantity
\[
 S_3(\rho)=R_{\rm opt}-L_{\rm disc}-L_{\rm unsat},                 \tag{SF5}
\]
where the two losses are precisely the two expectations in GH7.
The first-hit labels do not restrict legal exchanges to equal labels.
Both row-preserving and column-preserving exchanges stay available.
There is no assertion that compatible blockers are bad coefficients.

\section{An adaptive stopping event must be predictable}
Condition only on the original marks, matching and order, before
the tuple draws. In each full block, the selected source set is a
uniform four-subset of its $L$ vertices, independently between blocks.
The ordered roles and original bits retain their original laws.

Consider a fixed policy which queries membership of previously
unqueried base vertices, using only previous membership answers,
the conditioned data, and independent auxiliary randomness.
Designate a probe \emph{before} seeing its answer.
Immediately before such a probe in block $i$, let $e_i$ be the
number of exposed vertices and $r_i$ the number of selected answers.
Then, on every positive-probability transcript,
\[
 \Prb(\text{probe is selected}\mid\text{transcript})
                        =\frac{4-r_i}{L-e_i}.                    \tag{SF6}
\]
Indeed the still-hidden selected subset is uniform of size $4-r_i$
among $L-e_i$ vertices. Adaptive choices disclose no other information
under the stated filtration. Revealing a tuple, its outgoing targets,
or a label selected using the full history is not this filtration.

\begin{proposition}
Suppose $e_i\le L/2$ before every designated probe. The event that
a fixed policy emits $k$ designated probes and all are selected has
complete-original-law probability at most
\[
                         (8/L)^k.                              \tag{SF7}
\]
The bound allows intervening membership queries, failure-dependent
choices of blocks, and early stopping. It does not allow discarding
a failed designated probe and calling a later success the same probe.
\end{proposition}
\begin{proof}
SF6 bounds the conditional success probability by $8/L$ at each
designated probe. The event that a next probe is attempted is known
before its answer. Conditional expectation bounds the probability
of extending an all-success prefix by another factor $8/L$.
Induction proves SF7. There are finitely many distinct source queries,
so no infinite stopping argument is required. Average the original
marks, matching, order and independent auxiliary variables afterwards.
Remainder vertices cannot succeed and hubs are not counted as probes.
\end{proof}

Thus, for $r\le s$ policies fixed before sampling and normalized
first-removal charges $z_v=Z_v(U)/F(U)$,
\[
 \E\sum_v z_v1_{E_v}\le3\min(1,r(8/L)^k).                       \tag{SF8}
\]
This follows from GW9 and SF7, without independence between $z_v$
and $E_v$. It is a valid adaptive event budget, not a construction
of events dominating SF5.

\section{Where the attempted proof stops}
The missing assertion is a circuit-dependent domination of SF5 by
a joint dimension-independent defect charge and charges of the form
SF8, with enough probes and a proved aggregate coefficient budget.
SF2 supplies the coefficients; it does not supply that domination.

There are two distinct issues. First, replacing $\mu$ by $\nu$
in SF3 drops the displayed density factor. Second, the direct
selection of a completed pair, a maximizing Hall subset, or a
clipping threshold uses the full $(H,I)$, including the tuples.
The pins obtained that way are not designated before their membership
is known. SF6 therefore cannot be applied to that selection.
No alternative predictable exploration has been proved to retain
the required signed residual while keeping the exposure budget.
The small-size hypothesis alone has not supplied either missing
implication. These are proof gaps, not refutations of SF1, FH16,
GH15, or the residual-specific first-removal route.

If SF1 were established, FH9 and $a\le2^{-136}$ would give
\[
 |F_3(\rho)|<15e^{-N/2^{136}}\le15e^{-1024}<1/64.
\]
FH15 would then give $-\E_\Delta\sigma h_Q<17/64$, contradicting
BP3's lower bound above $7/16$. Hence $s\ge e^{aN}$.
With $D$ the fixed base-graph degree and
$D_*=\sum_{i=1}^{14L}D^i$, the existing variable convention gives
$M_n\le(D_*+6)N$ and therefore
$s\ge2^{aM_n/(D_*+6)}$. This is a conditional implication only.
FH12 still has factor
$1+1/[16(1+3/(5\cdot2^{65}))^q]$ and exponent $2^{-137}$;
FH13 still has exponent $2^{-136}$ at the old DS16 factor.
Neither benchmark has been improved here.

\paragraph{Exact replay and its scope.}
\path{first_hit_signed_check.py} checks the new first-hit partition
on a shared-square DAG and the measures SF3 on the reviewed two-block
factor fixture, retaining both optional directions and the original
dependent clipping. It does not rerun the previous Hall optimization.
A separate ten-vertex uniform-four-subset experiment has sixteen
membership atoms. After querying vertex zero, its two designated
probes are $(1,2)$ if zero was selected and $(2,3)$ otherwise.
Their conditional all-success probabilities are $1/12$ and $1/6$;
the unconditional probability is $2/15$. The inverse conditional
likelihood on successful paths has expectation one.
The inspected Jacobian conditional-expectation operation checks
this new finite law, not SF1 or a production-size enumeration.
Dedicated JSON files preserve its payload, typed output and execution
status. The all-size statements above are proved in prose.
\end{document}
